% ---------------------------------------------------------------------
%  Glossaire français-anglais
% ---------------------------------------------------------------------
\chapter{Glossaire}\label{ann:glossaire}

Les termes sont classés par ordre alphabétique. L'équivalent anglais, entre parenthèses, est celui
que l'on rencontre dans la documentation des outils.

\begin{description}[leftmargin=0pt,itemsep=3pt,font=\normalfont\bfseries]
\item[ACID] (\emph{ACID}) Propriétés d'une transaction : atomicité, cohérence, isolation et
  durabilité. Une transaction s'applique en entier ou pas du tout.
\item[Architecture en médaillon] (\emph{medallion architecture}) Organisation des données en trois
  couches successives : Bronze, Silver et Gold.
\item[Batch] (\emph{batch processing}) Voir \emph{traitement par lots}.
\item[Bronze, Silver, Gold] Couches de l'architecture en médaillon : données brutes, données
  propres et contrôlées, données métier prêtes pour l'analyse.
\item[Catalogue] (\emph{catalog}) Service qui associe le nom d'une table à ses métadonnées
  courantes, pour que tous les moteurs voient la même version.
\item[Checkpoint] (\emph{checkpoint}) Dossier où une requête en continu note les positions lues et
  écrites, pour reprendre sans perte ni doublon après un arrêt.
\item[Clé de partition] (\emph{partition key}) Dans Kafka, valeur qui détermine la partition d'un
  message ; dans une table, colonne qui détermine le rangement des fichiers.
\item[Compactage] (\emph{compaction}) Regroupement de nombreux petits fichiers d'une table en
  quelques gros fichiers.
\item[Dimension] (\emph{dimension}) Table qui décrit le contexte d'un fait : client, produit,
  vendeur, date.
\item[DAG] (\emph{Directed Acyclic Graph}) Graphe orienté sans cycle ; dans Airflow, description
  des tâches d'un pipeline et de leur ordre.
\item[ELT, ETL] (\emph{Extract, Load, Transform} ; \emph{Extract, Transform, Load}) Chargement puis
  transformation des données, ou l'inverse.
\item[Entrepôt de données] (\emph{data warehouse}) Base orientée analyse, alimentée par les
  applications, historisée et structurée.
\item[Fédération] (\emph{federated query}) Requête qui joint, en une fois, des données de plusieurs
  systèmes.
\item[Fenêtre de temps] (\emph{window}) Intervalle de temps sur lequel on agrège un flux, par
  exemple une journée.
\item[Format de table] (\emph{table format}) Couche de métadonnées (Iceberg, Delta Lake, Hudi) qui
  fait d'un ensemble de fichiers une table transactionnelle et versionnée.
\item[Fuite de données] (\emph{data leakage}) Utilisation, pour entraîner un modèle,
  d'informations qui ne seraient pas disponibles au moment de la prédiction.
\item[Groupe de consommateurs] (\emph{consumer group}) Ensemble de lecteurs Kafka qui partagent une
  même position de lecture.
\item[Heure de l'événement] (\emph{event time}) Moment où un événement s'est produit, par
  opposition à l'heure de son traitement (\emph{processing time}).
\item[Idempotence] (\emph{idempotence}) Propriété d'un traitement qui, relancé, ne change pas le
  résultat.
\item[Instantané] (\emph{snapshot}) État d'une table à un moment donné ; chaque écriture en crée un.
\item[Lac de données] (\emph{data lake}) Stockage de fichiers dans leur format d'origine, à faible
  coût, sans garantie transactionnelle.
\item[Lakehouse] Architecture qui ajoute à un lac de données les garanties d'un entrepôt, grâce à
  un format de table ouvert.
\item[Lineage] (\emph{data lineage}) Graphe qui relie chaque table aux traitements qui la
  produisent et aux tables qu'ils lisent.
\item[Micro-lot] (\emph{micro-batch}) Petit lot de données traité à chaque déclenchement d'une
  requête en continu.
\item[Offset] (\emph{offset}) Numéro d'un message dans sa partition Kafka.
\item[OLAP, OLTP] (\emph{OnLine Analytical / Transaction Processing}) Usages analytique et
  transactionnel des données.
\item[Orchestration] (\emph{orchestration}) Enchaînement automatique des tâches d'un pipeline, avec
  calendrier, relances et historique.
\item[Partition] (\emph{partition}) Dans Kafka, morceau d'un topic dans lequel l'ordre est garanti ;
  dans une table, sous-ensemble de fichiers qui partagent une valeur de colonne.
\item[Quarantaine] (\emph{quarantine}) Table où l'on range les lignes rejetées par un contrôle de
  qualité, avec la raison du rejet.
\item[Retard d'un groupe] (\emph{consumer lag}) Nombre de messages qu'un groupe de consommateurs n'a
  pas encore lus.
\item[Schéma en étoile, en constellation] (\emph{star schema}, \emph{galaxy schema}) Une table de
  faits entourée de ses dimensions ; plusieurs tables de faits qui partagent leurs dimensions.
\item[Stockage objet] (\emph{object storage}) Stockage de fichiers identifiés par un chemin, comme
  Amazon S3, accessible par une interface web.
\item[Streaming] (\emph{stream processing}) Voir \emph{traitement en continu}.
\item[Table de faits] (\emph{fact table}) Table qui enregistre des événements mesurables, avec
  leurs mesures et leurs clés vers les dimensions.
\item[Topic] (\emph{topic}) Dans Kafka, file de messages identifiée par un nom.
\item[Traitement en continu] (\emph{stream processing}) Traitement des événements dès leur arrivée.
\item[Traitement par lots] (\emph{batch processing}) Traitement, à heure fixe, des données
  accumulées.
\item[Vecteur] (\emph{embedding}) Suite de nombres qui résume le contenu d'un texte ou d'une image.
\item[Voyage dans le temps] (\emph{time travel}) Lecture d'une table telle qu'elle était à un
  instantané passé.
\item[Watermark] (\emph{watermark}) Limite au-delà de laquelle un événement en retard n'est plus
  attendu ; elle borne l'état d'une requête en continu.
\end{description}
